%% sections/results.tex — v15: expanded joint evaluation results

\section{Results}

\subsection*{Functional coverage}
Confluencia is the only platform evaluating circRNA-specific dimensions and providing cross-modality Go/Conditional/No-Go decisions (Table~\ref{tab:coverage}).

\subsection*{RNACTM dynamics}
Half-life consistency is by design ($<$4.1\% deviation from parametrization sources); constitutes consistency verification, not independent validation. The informative metric is the expression window deviation (40\,h vs 48\,h, 17\%). Endosomal escape 5.2\% slightly exceeds the 1--5\% range \citep{gilleron2013}. The 17\% deviation reflects the gap between isolated rate measurements and integrated compartment dynamics---clinically relevant for dosing interval decisions.

\subsection*{Prediction accuracy}
Epitope: WMA R$^2{=}0.82$, MAE 39\% improvement over Ridge ($N{=}300$). Drug: binding R$^2{=}0.95$, efficacy R$^2{=}0.60$ on 91K data; removing 2048-dim Morgan FP improves binding R$^2$ from 0.67 to 0.95, validating sample-size-adaptive design. MHC: allele-specific AUC${=}0.80$ on 52K binary data (246 alleles); focused-repertoire alleles match NetMHCpan (AUC${>}0.92$) (Table~\ref{tab:allele_auc}).

\begin{table}[t]
\centering
\caption{Per-allele MHC-I binding AUC on 52\,K IEDB binary data.}
\label{tab:allele_auc}
\begin{tabular}{lccc}
\toprule
\textbf{Allele} & \textbf{AUC} & \textbf{$n$} & \textbf{NetMHCpan range} \\
\midrule
A*33:03 & 0.95 & 315 & 0.92--0.96 \\
A*68:01 & 0.86 & 312 & 0.85--0.90 \\
A*24:02 & 0.70 & 604 & 0.85--0.90 \\
A*02:01 & 0.67 & 2144 & 0.92--0.96 \\
A*03:01 & 0.59 & 540 & 0.85--0.90 \\
\bottomrule
\multicolumn{4}{l}{\small Overall: AUC 0.83, 246 alleles.}
\end{tabular}
\end{table}

\subsection*{Evaluation stability}
Under $\pm$50\% weight perturbation: $\Psi$ 100\% Go, unmodified 100\% Conditional; Kinetics most critical (removal shifts $\Psi$ Go$\to$Conditional). Immunogenicity rankings agree with literature IFN data ($\rho{=}0.93$, direction consistency over 5 conditions).

\subsection*{Joint evaluation}
Real-data joint evaluation produces 5/5 valid dimensions with meaningful decision differentiation. Lapatinib (TNBC-targeted) yields composite${=}0.47$ (Conditional) vs aspirin composite${=}0.36$ (No-Go). Gene signature varies with target expression (high TROP2 0.54 vs low TROP2 0.33); dose changes affect kinetics (200\,mg 0.73 vs 10\,mg 0.31). The uncertainty-adaptive weighting $(1{-}u_i)^2$ auto-downweights under-constrained dimensions: Gene Signature ($u{=}0.50$) effective weight drops from 15\% to 3.75\%, preventing noisy features from dominating decisions.

\subsection*{Modification prioritization}
$\Psi$ is Go (composite 0.67); m6A and unmodified are Conditional (0.55/0.53) (Table~\ref{tab:casestudy}). The gap reflects $\Psi$'s dual benefit: extended half-life (15.61\,h) and low immunogenicity.

\begin{table}[t]
\centering
\caption{Modification prioritization for TNBC-targeting circRNA.}
\label{tab:casestudy}
\begin{tabular}{lcccc}
\toprule
\textbf{Modification} & \textbf{Half-life (h)} & \textbf{Immunogenicity} & \textbf{Composite} & \textbf{Decision} \\
\midrule
Unmodified & 6.24 & High & 0.53 & Conditional \\
$\Psi$ & 15.61 & Low & 0.67 & Go \\
m6A & 11.24 & Moderate & 0.55 & Conditional \\
\bottomrule
\end{tabular}
\end{table}
